In mathematics, a monomial is a constant multiplied by a variable raised to a power:

\[
    m \cdot X^{p}
\]

where:

\begin{itemize}

    \item \(m\) (\textbf{coefficient}) is a constant value.

    \item \(X\) (\textbf{variable}) is a variable value.

    \item \(p\) (\textbf{exponent}) is the power the variable is raised to.

\end{itemize}

Examples of monomials are \(2X^{5}\) and \(X^{-3}\). The goal of this exercise is to write a \texttt{Monomial} class in Java that represents this kind of data and performs the usual operations with it.

\subsection{Implementation}

Write a \texttt{Monomial} class with the following requirements:

\begin{itemize}

    \item \textbf{Attributes}:

        \begin{itemize}

            \item \texttt{double coefficient}: the coefficient of the monomial.

            \item \texttt{int exponent}: the power the variable is raised to.

            \item The variable is always \(X\), so it is not stored.

        \end{itemize}

    \item \textbf{Methods}:

        \begin{itemize}

            \item \textbf{Constructor}: creates a monomial from a coefficient and an exponent.

            \item \texttt{describe()}: returns a \texttt{String} with the written form of the monomial, made by joining the coefficient, \texttt{X\textasciicircum{}} and the exponent with \texttt{+}. For example, the monomial with coefficient \texttt{2} and exponent \texttt{3} is described as \texttt{"2.0X\textasciicircum{}3"} (a \texttt{double} is always written with a decimal point).

            \item \texttt{addition()}: given another \texttt{Monomial}, \textbf{returns a new monomial} with the sum of both. Two monomials only add up to a single monomial when they have the same exponent: the coefficients are added and the exponent stays the same. If the exponents differ, the method returns \texttt{null}.

            \item \texttt{scalarProduct()}: given a number, \textbf{returns a new monomial} with the product of the monomial by that number.

            \item \texttt{product()}: given another \texttt{Monomial}, \textbf{returns a new monomial} with the product of both: the coefficients are multiplied and the exponents are added.

            \item \texttt{evaluate()}: given a number, returns the value of the monomial when \(X\) takes that value. \texttt{Math.pow(a, b)} computes \(a^{b}\).

        \end{itemize}

\end{itemize}

None of the operations changes the monomials it works with: each one creates a new monomial with \texttt{new} and returns it.

\subsection{Recommendations}

\begin{itemize}

    \item \textbf{Incremental implementation}: write and test each method on its own before moving on to the next one, making sure each part works before building the rest of the class.

    \item \texttt{this}: write \texttt{this.} before every attribute, such as \texttt{this.coefficient}, even where no parameter has the same name, as in the slides. It makes clear which values belong to the monomial itself and which to \texttt{other}.

\end{itemize}

\subsection{Evaluation}

The \texttt{Main} class provided creates several monomials and prints the result of each operation next to the value it should have. To complete the exercise:

\begin{enumerate}

    \item \textbf{Class functionality}:

        \begin{itemize}

            \item Check that the constructor stores both values.

            \item Check the text returned by \texttt{describe()}.

        \end{itemize}

    \item \textbf{Monomial operations}:

        \begin{itemize}

            \item Check that additions and products give the right coefficient and exponent, and that the original monomials do not change.

        \end{itemize}

    \item \textbf{Clean code and documentation}:

        \begin{itemize}

            \item Make sure the code is clean and legible, with explanatory comments where necessary.

        \end{itemize}

\end{enumerate}
